\chapter{Profile-Guided, Learned, and Spatial Compilation}
\label{ch:profile-spatial-compilation}

A legal specialization can make a program slower.
Its guard costs time, its failure path may reconstruct state, and its code may displace another hot function from cache.
Moving work to an accelerator can also lose when transfer and launch overhead exceed the computation saved.
The compiler therefore needs a profitability model as well as a correctness argument.
Profiles and learned models supply evidence for the former; they do not replace the latter.

\section{Attach probabilities to observations, not truths}
\label{sec:ch68-probabilities}

A \term{profile annotation} records observations together with the region, version, and workload context that produced them.
A branch taken $90$ times in $100$ observations has an empirical frequency of $0.9$.
It is not proved to take that branch with probability $0.9$ on future inputs.
The observations may be correlated, the workload may change, and instrumentation may miss events.
Correctness must still cover every behavior permitted by the program contract.

For a speculative region, let $p$ be an estimated success probability, $g$ the guard cost, $f$ the successful fast-path cost, $d$ recovery cost, and $b$ baseline cost after failure.
Under this simplified mutually exclusive outcome model, the expected cost is
\[
C(p)=g+pf+(1-p)(d+b).
\]
This equation assumes the supplied costs already include the relevant work for each outcome.
It must not count the same prefix in both fast and fallback paths or ignore a repeated side effect.

Take $g=1$, $f=4$, $d=20$, and $b=10$ in synthetic cost units.
Then $C(p)=31-26p$, so the specialization beats baseline cost $10$ only when $p>21/26$, approximately $0.808$.
Table~\ref{tab:ch68-risk} computes two possible workloads.
No measured speedup is claimed: these numbers are an inspectable decision example.

\begin{table}[htbp]
\caption{The same correct specialization can win or lose as the success frequency changes.}
\label{tab:ch68-risk}
\centering
\begin{tabular}{lll}
\toprule
Estimated success & Expected speculative cost & Baseline comparison \\
\midrule
$0.9$ & $7.6$ & saves $2.4$ \\
$0.7$ & $12.8$ & loses $2.8$ \\
\bottomrule
\end{tabular}
\end{table}

\section{Make uncertainty affect policy}
\label{sec:ch68-uncertainty}

If plausible success probabilities lie in $[0.7,0.9]$, this candidate lacks a robust benefit under that interval model.
Because $C(p)$ decreases with $p$ here, its worst value over the interval is $C(0.7)=12.8$.
A policy that requires improvement throughout the interval rejects or postpones specialization.
The interval is an input assumption in this example, not a statistically calibrated confidence interval.

Algorithm~\ref{alg:ch68-policy} separates legality, estimated benefit, compilation cost, and a safety margin for profitability.
For a remaining invocation count $H$, compilation cost $K$, and per-call saving $s$, the simplest amortization condition is $Hs>K$.
Using the optimistic saving $2.4$ and $K=240$ requires more than $100$ remaining invocations.
If $H$ itself is uncertain, the policy must account for that uncertainty too.

\begin{algorithm}[htbp]
\caption{Choose a specialization using a conservative finite candidate model.}
\label{alg:ch68-policy}
\begin{pseudocode}
\Require candidates with established legality, cost intervals, horizon estimate, budget
\Ensure selected candidate satisfies the stated profitability policy
\State retain the baseline as the default choice
\ForAll{legal candidates within compilation and memory budgets}
  \State compute the worst modeled execution cost over the supplied uncertainty set
  \State include compilation, recovery, transfer, and retention costs as applicable
  \If{modeled total cost improves on baseline by the required margin}
    \State retain the best candidate under the declared objective
  \EndIf
\EndFor
\State \Return the selected version or baseline
\end{pseudocode}
\end{algorithm}

For $M$ candidates and an uncertainty set evaluated in $U$ scenarios, the direct policy costs $O(MU)$ model evaluations.
Analytic affine bounds reduce our interval example to endpoint evaluation.
The policy is only as reliable as its model and uncertainty assumptions; rejecting doubtful work is a performance choice, not a proof that the baseline is always fastest.

\section{Learn costs while keeping legality explicit}
\label{sec:ch68-learning}

A \term{learned cost model} predicts a property such as latency, throughput, code size, or relative candidate ranking from features and training observations.
The compiler can use the prediction to order search or choose among legal schedules.
Halide's learned autoscheduling work combines a learned model with search over schedules~\cite{adams-2019-learning}.
LLVM's MLGO infrastructure integrates learned decisions into optimization settings such as inlining and register-allocation policy~\cite{llvm-2026-mlgo}.

The smallest runnable example fits $\widehat{t}(n)=a+bn$ by least squares.
Its synthetic samples are $(1,7)$, $(2,10)$, $(3,13)$, and $(4,16)$.
The sums are $\sum n=10$, $\sum t=46$, $\sum n^2=30$, and $\sum nt=130$.
With four samples, the denominator is $4\cdot30-10^2=20$, giving $b=(4\cdot130-10\cdot46)/20=3$ and $a=(46-3\cdot10)/4=4$.
The predicted value at $n=5$ is $19$.
This checks fitting arithmetic, not generalization to real hardware.

Figure~\ref{fig:ch68-learning} shows the model's position in the optimization loop.
Training labels require controlled measurement, a specified objective, and separation between training and evaluation workloads.
Feature distributions can shift across programs, machines, and compiler versions.
For a candidate unlike the training data, a fallback heuristic or additional measurement may be more reliable than a confident extrapolation.

\begin{figure}[htbp]
\centering
\begin{tikzpicture}[node distance=14mm]
\node[basic block] (legal) {legal candidate schedules};
\node[basic block,below=of legal] (model) {features and cost prediction};
\node[basic block,below=of model] (select) {search and candidate selection};
\node[basic block,below=of select] (measure) {validate and measure output};
\draw[cfg edge] (legal)--(model); \draw[cfg edge] (model)--(select); \draw[cfg edge] (select)--(measure);
\end{tikzpicture}
\caption{Learning influences candidate selection while semantic constraints remain an independent gate.}
\label{fig:ch68-learning}
\end{figure}

\section{Represent temporal relationships with a stated meaning}
\label{sec:ch68-temporal}

Two functions that execute close together may benefit from nearby code placement even if neither calls the other directly.
We can record a \term{temporal affinity graph} whose nodes are artifacts and whose weighted edges count occurrences within a specified observation window.
The window definition, directionality, normalization, and version identities are part of the data schema.
The phrase describes this chapter's model; it is not an attribution of a universal standard graph representation.

Temporal affinity is different from a happens-before relation, a dependence edge, or a proof of legal instruction reordering.
The observed sequence $A,B,A$ yields proximity evidence for $A$ and $B$ but does not authorize exchanging their execution.
Similarly, an IR annotated with branch frequencies must still preserve behavior for every permitted input; the annotations do not turn it into a probabilistic program.
Probabilistic information about execution and semantics of probabilistic programs are separate concepts.

\section{Place related artifacts under explicit constraints}
\label{sec:ch68-placement}

Consider three equal-sized artifacts in consecutive unit slots.
Their symmetric pair weights are $w_{AB}=9$, $w_{BC}=8$, and $w_{AC}=1$.
For a permutation $\pi$, use the toy distance objective
\[
J(\pi)=\sum_{i<j}w_{ij}|\pi(i)-\pi(j)|.
\]
Table~\ref{tab:ch68-placement} evaluates three orders; reversal preserves their costs.
Placing $B$ between $A$ and $C$ gives $9+8+2=19$.
The other two middle choices yield $26$ and $27$.

\begin{table}[htbp]
\caption{The toy placement objective favors short distances on heavily weighted pairs.}
\label{tab:ch68-placement}
\centering
\begin{tabular}{lll}
\toprule
Order & Weighted distance calculation & Total \\
\midrule
$A,B,C$ & $9+8+2$ & $19$ \\
$A,C,B$ & $18+8+1$ & $27$ \\
$B,A,C$ & $9+16+1$ & $26$ \\
\bottomrule
\end{tabular}
\end{table}

Algorithm~\ref{alg:ch68-place} is exact because it enumerates every permutation.
It costs $O(n!n^2)$ time for $n$ artifacts with dense pair weights, so it is suitable only for tiny instances and checking heuristics.
Large layout systems use heuristics under real alignment, section, branch-range, and code-size constraints.
Byte distance alone does not model cache sets, prefetching, page crossings, or branch predictors.

\begin{algorithm}[htbp]
\caption{Find the exact optimum of the small unit-slot placement model.}
\label{alg:ch68-place}
\begin{pseudocode}
\Require $n$ equal-sized artifacts, symmetric nonnegative pair weights
\Ensure a permutation minimizing the stated weighted-distance objective
\State $best\gets$ infinity; $order\gets$ none
\ForAll{permutations $p$ of the artifact IDs}
  \State construct the position of each artifact in $p$
  \State $cost\gets\sum_{i<j}w_{ij}|p^{-1}(i)-p^{-1}(j)|$
  \If{$cost<best$}
    \State $best\gets cost$; $order\gets p$
  \EndIf
\EndFor
\State \Return $best,order$
\end{pseudocode}
\end{algorithm}

Figure~\ref{fig:ch68-placement} displays the winning order with its two strongest affinities adjacent.
The same idea can inform code layout, object clustering, or device placement, but each application needs its own capacity and access model.
Rearranging physical objects also requires updating or preserving every reference according to the runtime's identity contract.

\begin{figure}[htbp]
\centering
\begin{tikzpicture}[node distance=23mm]
\node[basic block] (a) {slot 0: A};
\node[basic block,right=of a] (b) {slot 1: B};
\node[basic block,right=of b] (c) {slot 2: C};
\draw[cfg edge] (a)--node[above,font=\small] {$9$}(b);
\draw[cfg edge] (b)--node[above,font=\small] {$8$}(c);
\draw[dashed] (a.south) to[bend right=25] node[below,font=\small] {$1$, distance $2$}(c.south);
\end{tikzpicture}
\caption{The displayed arrows indicate affinity weights, not control-flow or ordering permissions.}
\label{fig:ch68-placement}
\end{figure}

\section{Include transfer costs in heterogeneous execution}
\label{sec:ch68-heterogeneous}

Suppose a CPU takes $4n$ units for a task while a device takes $120+n$, including launch and transfer in the fixed term.
The break-even point solves $4n=120+n$, hence $n=40$.
At $n=41$ the device costs $161$ versus the CPU's $164$.
Selecting the device solely because its kernel coefficient is smaller would lose on small inputs.

Real placement can depend on data residency, queue delay, bandwidth, device memory capacity, energy, and synchronization.
A result consumed immediately on the CPU may require a return transfer; a pipeline of device consumers may amortize that cost.
Different devices may also implement different floating-point or atomic behavior, so semantic compatibility is a prerequisite to cost comparison.
A runtime-owned optimizer can update placement as conditions change, but needs hysteresis to avoid repeatedly paying migration costs near a decision boundary.

\section{Exercises}
\label{sec:ch68-exercises}

\begin{enumerate}
\item\label{ex:ch68-threshold} Derive the exact success threshold for the speculative-cost example and compute both endpoint costs for $p\in[0.7,0.9]$.
\item\label{ex:ch68-fit} Recompute the linear fit from the four supplied samples and predict cost at $n=6$.
State what this calculation does not establish about measured hardware behavior.
\item\label{ex:ch68-order} Enumerate the six three-artifact orders, group reversals, and verify the optimum and both other costs.
\item\label{ex:ch68-device} Solve the CPU/device break-even point and compare costs at $n=20$ and $n=80$.
Explain how retaining data on the device can change the model.
\end{enumerate}

\section{Further reading}
\label{sec:ch68-reading}

The Halide autoscheduler paper provides a concrete learned-model and search design~\cite{adams-2019-learning}.
LLVM's MLGO documentation explains integration of learned optimization decisions into compiler infrastructure~\cite{llvm-2026-mlgo}.
The companion implements the cost arithmetic, least-squares fit, and exhaustive tiny placement model with synthetic inputs; it does not report benchmark improvements.
